\section{Discussion}

We set out to let an agent choose how far to deviate from a safe policy without the act of choosing biasing the estimate that justified the choice. The resolution exploits the fact that the only shift between the safe policy's data and any candidate's deployment is the agent's own: it does not know how the world will respond, but it knows exactly how each candidate would distribute its actions. Conformal policy control parameterizes that choice as a likelihood-ratio threshold and calibrates it directly on the safe policy's data, so that selection is the calibrated act rather than a step that follows it. Because the threshold bounds how far the candidate can deviate from the safe policy, the resulting guarantee applies to the chosen policy itself, not merely each candidate in isolation.

More broadly, CPC instantiates a prescriptive use of conformal methods: rather than characterizing uncertainty under a fixed distribution, it searches over distributions to find one that satisfies a finite-sample guarantee. This inversion applies wherever one selects a data-generating process subject to a risk constraint, from experimental design, to active data collection, to generative model deployment. We hope that the ideas in this paper encourage further development of conformal methods as tools for decision-making, beyond their traditional role in uncertainty quantification.

The guarantees we provide are marginal: risk is controlled on average over context and calibration data draws.
This is the natural guarantee for deployment decisions, where one asks whether a policy is safe to deploy across a population.
It is less informative for individual decisions, where one asks whether a particular action is safe for a particular context.
Stronger conditional guarantees are possible but require additional assumptions or more conservative bounds \citep{thomas2019preventing}.
Our guarantees also depend on the stability of the context distribution.
If the context distribution begins to shift then recalibration may be necessary, or some risk control gap may need to be tolerated \citep{barber2023}.
Monitoring methods based on conformal martingales offer a principled approach to detecting such shifts \citep{prinster2025watch}.
We have also assumed the policy likelihoods are computable in closed form. If the policy likelihoods are intractable, then density ratio estimation \citep{sugiyama2012density} or neural ratio estimation \citep{cranmer2020frontier} techniques may be applicable.

A recurring theme in this work is that assumptions should reflect what the practitioner knows rather than what is theoretically convenient. CPC requires a safe policy, but any baseline whose risk the practitioner is willing to accept will do. It requires smoothness, but of the agent's own policy likelihood ratios, not of the unknown constraint loss function. It requires a risk tolerance, but as a direct declaration rather than an indirect penalty weight. Safety and exploration, in this view, need not run counter to each other. In the right balance they correct each other's deficiencies, avoiding complacency on the one hand and dissipation on the other. The balance must be set from what is actually known, not what is hoped. It can be enough, it turns out, to know what works, and to know what we want to do next.

% A recurring theme in this work is that assumptions should reflect what the practitioner knows rather than what is theoretically convenient. CPC requires a safe policy, but any baseline whose risk the practitioner is willing to accept will do. It requires smoothness, but of the agent's own policy likelihood ratios, not of the unknown constraint loss function. It requires a risk tolerance, but as a direct declaration rather than an indirect penalty weight. Each of these choices reflects the same principle: build the method around the practitioner's actual knowledge, and let the guarantees follow from that.

% Safety and exploration need not run counter to each other.
% In the right balance, they correct each other's deficiencies, avoiding complacency on the one hand and dissipation on the other. 
% However, it must be possible to determine an acceptable balance from what we actually know, not what we hope to be true.
% It can be enough, it turns out, to know what works, and to know what we want to do next.

